\scanpage{146}
\mspage{- 156 -}
\hfill n\textsuperscript{o} 280

\uline{Corollaire} - (Hilbert)- \uline{Soient} $k$ \uline{un corps}, $K$ \uline{une extension algébriquement close de} $k$, $P$ \uline{un} \uline{élément de} $B=k\,[X_1,\ldots,X_n]$, \uline{et} $n$ \uline{un idéal de} $B$. \uline{Les deux propriétés suivantes sont équivalentes} :

a) \uline{il existe un entier} $r$ \uline{tel que} $P^r\in n$ ;

b) \uline{pour tout $k$-homomorphisme} $f$ \uline{de} $B$ \uline{dans} $K$ \uline{tel que} $f(n)=\{0\}$, \uline{on a} $f(P)=0$.

En effet, $B$ est un anneau de Jacobson (Th.1) . Donc les polynomes de $B$ dont une puissance est dans $n$ sont ceux qui appartiennent à tout idéal maximal de $B$ contenant $n$ . Il suffit alors d'appliquer la prop.1c).

% archive sheet number: 145

